% Noyau d'hélium : deux protons et deux neutrons, distance entre les centres des protons (énoncé du défi D11)
\begin{tikzpicture}[x=1cm,y=1cm]
\shade[ball color=otMuted!45] (0.95,0.75) circle (0.62);
\shade[ball color=otMuted!45] (1.05,-0.75) circle (0.62);
\shade[ball color=otRed!55] (0,0) circle (0.62); \node[font=\otSemi\fontsize{15}{15}\selectfont,text=white] at (0,0) {$+$};
\shade[ball color=otRed!55] (2,0) circle (0.62); \node[font=\otSemi\fontsize{15}{15}\selectfont,text=white] at (2,0) {$+$};
\draw[otMuted,line width=0.4pt,dash pattern=on 1.5pt off 1.2pt] (0,0) -- (0,-1.55) (2,0) -- (2,-1.55);
\draw[otInk!70,line width=0.4pt,{Stealth[length=1.4mm]}-{Stealth[length=1.4mm]}] (0,-1.5) -- (2,-1.5) node[midway,below,pcLegende] {$d=\qty{2,0e-15}{m}$};
\draw[otMuted,line width=0.4pt] (2.45,0.42) -- (3.3,0.9); \node[pcLegende,anchor=west] at (3.35,0.9) {proton, charge $+e$};
\draw[otMuted,line width=0.4pt] (1.5,1.1) -- (3.3,1.6); \node[pcLegende,anchor=west] at (3.35,1.6) {neutron, charge nulle};
\end{tikzpicture}
\par\smallskip\legende{Modèle simplifié du noyau d'hélium 4. Les quatre nucléons sont serrés les uns contre les autres.}
